\documentclass[10pt,a4paper]{amsart}
\usepackage[margin=19mm]{geometry}
\usepackage{amsmath,amssymb,mathtools}
\usepackage[scr=rsfs,cal=euler]{mathalfa}
\usepackage{xfrac}
\usepackage[unicode,hidelinks,bookmarks=false]{hyperref}
\newcommand{\ie}{i.e.\ }
\newcommand{\cf}{cf.\ }
\newcommand{\resp}{respectively\ }
\newcommand{\Subsection}{\subsection}
\providecommand{\nobreakdash}{\nobreak\hbox{-}}
\setlength{\emergencystretch}{2em}
\multlinegap0pt
\usepackage{bm}
\usepackage{bbm}
\usepackage{dsfont}
\usepackage{xspace}
\usepackage{cancel}
\usepackage{mathtools}
\usepackage{enumitem}
\usepackage{amsthm}
\makeatletter
\@ifundefined{lemma}{\newtheorem{lemma}{Lemma}}{}
\@ifundefined{theorem}{\newtheorem{theorem}[lemma]{Theorem}}{}
\makeatother
\newcommand\setsuchas[2]{\left\{{#1}:{#2}\right\}}
\newcommand{\NNp}{{\mathbb{P}}}
\newcommand{\ZZ}{{\mathbb{Z}}}
\newcommand{\divides}[2]{#1 \mid #2}
\newcommand{\modd}[2]{#1\ \textup{(mod }#2\textup{)}}
\newcommand{\set}[1]{\left\{{#1}\right\}}
\newcommand{\ssift}[1]{\mathbf{s}({#1})}
\newcommand{\val}{\operatorname{val}}
\title[Greedy Regular Convolutions]{Greedy Regular Convolutions}
\author{Jan Snellman}
\date{Source: arXiv:2504.02795v4 (CC BY 4.0)}
\begin{document}
\maketitle

\noindent\emph{Source and licence.} Greedy Regular Convolutions. Version arXiv:2504.02795v4 (CC BY 4.0), distributed under the Creative Commons Attribution 4.0 International licence, as recorded in the arXiv licence metadata for this exact version and shown on the abstract page https://arxiv.org/abs/2504.02795v4. Full rights notice: licenses/arxiv-2504.02795-CC-BY-4.0.txt.\par\smallskip

\noindent\textit{Editorial scope.} Counted unit: the Theorem whose source label is \texttt{org2dba548} in the article (source0.tex lines 1267--1273, source label \texttt{org2dba548}), delivered together with the article's own complete original proof of it (source0.tex lines 1274--1348, one \texttt{proof} environment of 75 lines). No mathematics is written, repaired or re-derived: statement and proof are the article's own text line for line. Required context retained uncounted: none. Every definition, display and label of those lines is retained unchanged. Omitted from this package and not redistributed: 1--1266, 1349--1588. Nothing inside a retained excerpt is omitted or altered: every retained body is delivered exactly as the article's lines are written, so no omission receipt of any kind accompanies this package. Citations: the retained excerpts carry no citation command at all, so no bibliography entry of the article is transcribed and no reference is redistributed. Reference adaptation: none was needed and none was made; every reference inside the delivered unit resolves inside it, so no unresolved reference or citation is produced. Printed slips kept literal: no printed word, symbol, hypothesis or number of the counted statement or of its proof was altered. Layout and class: the article's own class is \texttt{article} with packages including inputenc, fontenc, graphicx, longtable, wrapfig, rotating, ulem, amsmath, amssymb, capt-of, hyperref, amsthm, enumitem, tabu; the wrapper is the lane's pinned \texttt{amsart} editorial wrapper at 10pt on A4, which restates the theorem-like environments the retained text needs and adds the article's own macro definitions that the retained text needs (\texttt{\textbackslash NNp}, \texttt{\textbackslash ZZ}, \texttt{\textbackslash divides}, \texttt{\textbackslash modd}, \texttt{\textbackslash set}, \texttt{\textbackslash setsuchas}, \texttt{\textbackslash ssift}, \texttt{\textbackslash val}), with extra wrapper packages (bm, bbm, dsfont, xspace, cancel, mathtools, enumitem). The delivered numbering is the unit's own. Assets: no retained span contains a figure environment, a graphics-inclusion command, a PDF-page inclusion, a table or a TikZ picture; no image file is copied into this package, the package declares no asset dependency and no image file is redistributed with it. Licence: version arXiv:2504.02795v4 of this article is distributed under the Creative Commons Attribution 4.0 International licence, as recorded in the arXiv licence metadata for this exact version and shown on the abstract page https://arxiv.org/abs/2504.02795v4. A full rights notice is delivered at licenses/arxiv-2504.02795-CC-BY-4.0.txt. Characters: every retained excerpt is pure ASCII, so the delivered engine, XeLaTeX with its default Latin Modern text font, reports no missing character.
\begin{theorem}
If \(p\) is a prime, \(k\) a positive integer, \(S=\set{p,p^{k}}\), \(M=\ssift{S}\) then if \(k\) is odd
\[M = \setsuchas{n \in \NNp}{\val_{p}(n) \equiv \modd{0}{2}},\]
and if \(k\) is even then
\[M = \setsuchas{n \in \NNp}{\val_{p}(n) \equiv \modd{0,2,4, \dots, k-2}{k+1}}.\]
\label{org2dba548}
\end{theorem}

\begin{proof}
Assume first that \(k\) is odd.
We want to prove that
\[\ssift{\set{p,p^{k}}} = \setsuchas{n \in \NNp}{\val_{p}(n) \equiv \modd{0}{2}}.\]

Write \(n=p^{a}r\) with \(\gcd(p,r)=1\). We proceed by induction on \(k\).

If \(a=0\) then neither \(p\) nor \(p^{k}\) divides \(n\), so \(n \in \ssift{\set{p,p^{k}}}\).

If \(a>0\) then \(n/p=p^{a-1}r\).
If \(a\) is even then \(a-1\) is odd,
so \(n/p \not \in \ssift{\set{p,p^{k}}}\), by the induction hypothesis.
If \(a\) is odd then \(a-1\) is even,
so \(n/p \in \ssift{\set{p,p^{k}}}\).

If \(a \ge k\) then \(n/p^{k}=p^{a-k}r\).
If \(a\) is even then \(a-k\) is odd,
so \(n/p^{k} \not \in \ssift{\set{p,p^{k}}}\), by the induction hypothesis.
If \(a\) is odd then \(a-k\) is even,
so \(n/p^{k} \in \ssift{\set{p,p^{k}}}\).

Thus for all \(d \in \set{p,p^{k}}\) we see that if \(\divides{d}{n}\) then
\(n/d \in \ssift{\set{p,p^{k}}}\) if and only if \(a\) is odd. Hence
\(n=p^{a}r \in \ssift{\set{p,p^{k}}}\) if and only if \(a\) is even, as stated.

Next assume that \(k>0\) is even. Write
\[n=p^{a}r, \qquad \gcd(p,r)=1,\]
and write
\[a= c(k+1) + d,  \qquad 0 \le d < k+1.\]
Then \(n/p = p^{a-1}r\) (if \(a>0\)) and \(n/p^{k} = p^{a-k}r\) (if \(a \ge k\)).
Write the set of congruence classes modulo \(k+1\) as a disjoint union
\begin{displaymath}
\frac{\ZZ}{(k+1)\ZZ}  = A \sqcup B
\end{displaymath}
with
\begin{align*}
A & =  \set{[0]_{k+1}, [2]_{k+1}, [4]_{k+1}, \dots, [k-2]_{{k+1}}} \\
B &=  \set{[1]_{k+1}, [3]_{k+1}, [5]_{k+1}, \dots, [k-1]_{k+1}, [k]_{k+1}}
\end{align*}
We want to prove that
\[n=p^{a}r \in \ssift{p,p^{k}} \qquad \iff \qquad [a]_{k+1} \in A.\]

Observe that
\begin{itemize}[label={\(\star\)}]
\item Subtracting \([1]_{k+1}\) from an element of \(A\) yields an element of \(B\).
\item Subtracting \([k]_{{k+1}}\) from an element of \(A\) yields an element of \(B\).
\item Subtracting \([k]_{k+1}\) from an element of \(B\) yields an element of \(A\).
\item Subtracting \([1]_{k+1}\) from an element of \(B \setminus \set{[k]_{k+1}}\) yields an element of \(A\);
subtracting \([k]_{k+1}\) from \([k]_{k+1}\) yields \([0]_{k+1} \in A\).
\end{itemize}

With a little book-keeping we see that if
\[n=p^{a}r, \qquad [a]_{k+1} \in A\]
then if \(1 \le a\) then
\[n/p = p^{a-1}r, \qquad  [a-1]_{k+1} \in B,\]
hence
\[n/p \not \in \ssift{\set{p,p^{k}}}.\]
Similarly,  if
\(a \ge k\) then
\[n/p^{k} = p^{a-k}r, \qquad [a-k]_{k+1} \in B,\]
hence
\[n/p^{k} \not \in \ssift{\set{p,p^{k}}}.\]
In both cases, there is no allowed factor in \(\ssift{p,p^{k}}\), so
\(n \in \ssift{p,p^{k}}\).

Conversely,
if \[n=p^{a}r,\qquad [a]_{k+1} \in B,\] then
if \(1 \le a\) with \(a \not \in [k]_{k+1}\) then
\[n/p = p^{a-1}r,\qquad [a-1]_{k+1} \in A,\]
hence
\[n/p  \in \ssift{p,p^{k}} \quad  \text{ so } \quad n  \not \in \ssift{p,p^{k}}.\]
Finally, if \(a \ge k\) (this includes the case \(a  \in [k]_{k+1}\))
then \[n/p^{k} = p^{a-k}r,\qquad [a-k]_{k+1} \in A,\] hence
\[n/p^{k} \in \ssift{p,p^{k}} \quad \text{ so } \quad n  \not \in \ssift{p,p^{k}}.\]
\end{proof}

\end{document}
